\section{相关技术综述}

本章系统梳理与本研究直接相关的四个技术领域：知识图谱与图结构检索增强生成、多模态文档理解与知识抽取、多智能体系统与工作流编排，以及专业领域智能体应用，分析各方向的技术演进脉络及其与本研究的关联，为后续方法设计奠定理论基础。

\subsection{知识图谱与图结构检索增强生成}

\subsubsection{传统RAG方法及其局限性}

检索增强生成（RAG）由Lewis等人\cite{lewis2020retrieval}于2020年提出，通过在推理阶段动态检索外部文档并将其作为上下文输入LLM，显著降低了模型幻觉问题。早期RAG系统以稠密向量检索（Dense Passage Retrieval，DPR）为核心，利用FAISS等近似最近邻算法在向量空间中检索语义相近段落。

然而，向量检索RAG存在固有局限：其一，检索粒度以独立文本块（chunk）为单位，缺乏跨文本的关系推理能力；其二，相似度度量仅捕捉浅层语义，难以处理多跳推理问题；其三，孤立段落的上下文丢失导致检索结果碎片化，影响生成质量。Self-RAG\cite{asai2023selfrag}引入自反思机制以动态决定是否检索，CRAG\cite{yan2024crag}提出上下文自适应压缩策略，Adaptive RAG\cite{jeong2024adaptive}根据问题复杂度自适应切换检索策略，上述工作均在向量检索框架内进行改进，未能从根本上解决多跳推理能力的不足。

\subsubsection{基于知识图谱的结构化检索}

知识图谱（Knowledge Graph，KG）为结构化知识存储和关系推理提供了天然载体。知识图谱技术综述\cite{xu2016kgreview}和知识表示学习研究\cite{liu2016krl}表明，实体、关系和语义表示的统一建模是知识推理与智能问答的重要基础。知识图谱与RAG的结合（KGRAG）\cite{kgrag2023}通过实体链接将检索转化为子图查询，在关系推理场景下显著优于向量检索。KGRAG方法首先识别问题中的实体，随后在图中定位对应节点，通过图遍历抽取相关子图作为上下文，实现了从孤立文档到结构化知识的检索范式转换。

微软研究院提出的GraphRAG\cite{edge2024graphrag}进一步引入社区检测和社区摘要（Community Summary）机制，通过Leiden算法划分图社区，并为每个社区生成摘要，实现从局部到全局的层次化知识表示。GraphRAG在长文档摘要任务上展现出超越传统向量检索的性能，尤其适合需要全局理解的问答场景。

HippoRAG\cite{gutierrez2024hipporag}借鉴人类海马体记忆机制，构建个性化知识图谱支持长期记忆检索；RAPTOR\cite{sarthi2024raptor}通过递归文本摘要树构建多层次上下文索引。上述方法均展示了图结构在知识组织和多跳检索方面的优势。

\subsubsection{现有图RAG方法的不足}

尽管图结构RAG在文本知识检索上取得了重要进展，现有方法仍存在以下不足：第一，检索策略多为全局社区摘要或固定跳数遍历，缺乏以语义向量检索为种子的自适应子图扩展机制；第二，上述方法均针对纯文本信息，无法处理图像、流程图等多模态信息；第三，检索结果的格式化方式缺乏面向下游LLM推理的结构化设计。本研究提出的GraphRAG方法针对上述不足进行了系统性改进，将在第三章详述。

\subsection{多模态文档理解与知识抽取}

\subsubsection{视觉语言模型与多模态理解}

随着视觉语言模型（Vision-Language Model，VLM）的快速发展，文档图像的自动理解成为可能。CLIP\cite{radford2021clip}通过对比学习建立图文语义对齐，为多模态表示学习奠定了基础。GPT-4V\cite{openai2023gpt4}、LLaVA\cite{liu2023llava}等多模态大模型进一步实现了对复杂图表、流程图和表格的理解与描述能力。

在文档解析方面，DocLayoutYOLO\cite{doclayoutyolo2024}和TableTransformer\cite{smock2022tabletransformer}等专用模型分别针对文档版式检测和表格结构解析进行了优化。这些工作共同推动了多模态文档理解从简单OCR向深度语义理解的演进。

\subsubsection{多模态实体抽取与图谱构建}

传统知识图谱构建依赖于对纯文本的命名实体识别（NER）和关系抽取（RE），对图像和图形中的结构化信息无能为力。近年来，研究者开始探索多模态实体抽取方法。MMC\cite{liu2023mmc}通过大规模图表指令数据提升模型对图表的跨模态理解能力；在电力故障诊断场景中，已有研究将多模态知识图谱与视觉检测模型结合，用于变电站设备故障的自主预警和诊断\cite{xiao2022substationkg}。

在电力系统领域，单线图（Single Line Diagram，SLD）和网络拓扑图是重要的知识载体，蕴含丰富的设备连接关系。然而，现有方法尚未能有效解析电力系统专有图纸，将其中的实体（变压器、线路、断路器等）和关系（连接、保护等）自动转化为图谱节点和边。本研究利用VLM对电力图像进行描述，再通过LLM进行实体-关系抽取，实现了电力多模态文档的统一知识图谱构建。

\subsubsection{嵌入模型与多语言语义表示}

知识图谱的向量化检索依赖高质量的文本嵌入模型。BGE-M3\cite{chen2024bgem3}由智源研究院提出，支持超过100种语言的语义嵌入，在单向量检索、稀疏检索和多向量交互三种检索模式上均取得了优异性能。BGE-M3的多粒度表示能力使其在中英文混合的电力领域文档中具有显著优势，本研究选用BGE-M3生成1024维语义嵌入，并通过Neo4j向量索引实现高效的近似最近邻检索。

\subsection{多智能体系统与工作流编排}

\subsubsection{基于LLM的多智能体框架}

多智能体系统（Multi-Agent System，MAS）的研究最早源于分布式人工智能领域，近年随LLM能力的突破而获得新的发展动力。AutoGen\cite{wu2023autogen}设计了以可对话智能体为核心的编程接口，通过智能体间的多轮对话协作完成代码生成、数据分析等复杂任务，并支持人机混合协作模式。

MetaGPT\cite{hong2023metagpt}将软件工程中的角色分工引入多智能体系统，通过标准化输入输出协议规范智能体间通信，定义了产品经理、架构师、工程师等专业角色，实现了端到端的软件开发自动化。CAMEL\cite{li2023camel}探索了基于角色扮演的多智能体对话范式，通过提示工程引导智能体扮演特定角色。CrewAI\cite{crewai2024}以"船员"隐喻组织智能体团队，支持任务的层级分解和并行执行。

\subsubsection{工作流编排框架}

LangGraph\cite{langgraph2024}基于有向图（Directed Graph）的计算模型，将智能体工作流抽象为节点（Node）和边（Edge），支持条件分支、循环结构和并行子图的精确控制，并提供状态持久化和检查点机制。TaskWeaver\cite{taskweaver2024}以代码片段为执行单元的智能体编排机制，通过将用户请求转化为Python代码任务提升执行可靠性。

上述框架均依赖人工预先设计工作流拓扑，需要开发者以编程方式指定智能体间的调用关系。这种方式灵活性高，但难以应对开放域的动态任务需求——不同用户请求可能需要完全不同的工作流结构，人工预设无法覆盖所有场景。

\subsubsection{工作流自动生成方法}

工作流自动生成旨在根据任务描述，无需人工干预地生成最优智能体协作方案。现有研究主要依赖基于LLM的直接生成，通过提示工程或思维链机制引导LLM输出工作流结构\cite{wei2022chain}，但生成质量依赖模型内部知识，缺乏对外部领域知识和智能体能力约束的显式利用。

在智能体匹配方面，现有框架大多依赖LLM隐式的角色分配能力，缺乏对智能体能力（Capability）和工具（Tool）的显式建模与量化匹配。本研究提出的子任务-智能体匹配算法（算法C），通过三维评分函数系统化解决角色分配准确率问题，将在第四章详细介绍。

\subsection{本章小结}

本章从检索增强生成、多模态文档理解和多智能体框架三个维度梳理了相关技术的发展现状。在图结构RAG方面，GraphRAG等方法在文本知识检索上取得了重要进展，但缺乏多模态处理能力和面向工作流生成的子图自适应扩展机制；在多模态理解方面，VLM和BGE-M3为多模态知识图谱构建提供了技术基础；在多智能体编排方面，LangGraph等框架提供了图计算支撑，但自动工作流生成和系统化智能体匹配仍是待解难题。本研究正是在上述基础上，提出涵盖知识图谱构建、GraphRAG检索、智能体匹配和工作流编排四个环节的完整方法体系，各核心算法将在后续章节详细阐述。
